% !Mode:: "TeX:UTF-8"

\hitsetup{
  %******************************
  % 注意：
  %   1. 配置里面不要出现空行
  %   2. 不需要的配置信息可以删除
  %******************************
  %
  %=====
  % 秘级
  %=====
  statesecrets={},
  natclassifiedindex={TM301.2},
  intclassifiedindex={62-5},
  %
  %=========
  % 中文信息
  %=========
  ctitleone={面向推理过程的},%本科生封面使用
  ctitletwo={大语言模型错误识别与改进优化},%本科生封面使用
  ctitlecover={面向推理过程的大语言模型错误识别与改进优化},%放在封面中使用，自由断行
  ctitle={面向推理过程的大语言模型错误识别与改进优化},%放在原创性声明中使用
  % csubtitle={一条副标题}, %一般情况没有，可以注释掉
  cxueke={工学},
  csubject={人工智能},
  caffil={计算学部},
  cauthor={许锐一},
  csupervisor={冯骁骋教授},
  % cassosupervisor={某某某教授}, % 副指导老师
  % ccosupervisor={某某某教授}, % 联合指导老师
  % 如果是第一封面的日期要手动设置，需要取消注释下一行，并将内容改为“规范”中要求的封面第一页最下方的日期
  % firstpagecdate={2022年6月},
  % 日期自动使用当前时间，若需指定按如下方式修改：
  cdate={2026 年 4 月 25 日},
  cstudentid={2022111446},
  cstudenttype={学术学位论文}, %非全日制教育申请学位者
  cnumber={no2022111446}, %编号
  % cpositionname={哈铁西站}, %博士后站名称
  cfinishdate={20XX年X月---20XX年X月}, %到站日期
  % csubmitdate={20XX年X月}, %出站日期
  % cstartdate={3050年9月10日}, %到站日期
  % cenddate={3090年10月10日}, %出站日期
  %（同等学力人员）、（工程硕士）、（工商管理硕士）、
  %（高级管理人员工商管理硕士）、（公共管理硕士）、（中职教师）、（高校教师）等
  %
  %
  %=========
  % 英文信息
  %=========
  etitle={Error Identification and Optimization for Large Language Models in Reasoning Processes},
  esubtitle={This is the sub title},
  exueke={Engineering},
  esubject={Computer Science and Technology},
  eaffil={\emultiline[t]{School of Mechatronics Engineering \\ Mechatronics Engineering}},
  eauthor={Yu Dongmei},
  esupervisor={Professor XXX},
  eassosupervisor={XXX},
  % 日期自动生成，若需指定按如下方式修改：
  edate={December, 2017},
  estudenttype={Master of Art},
  %
  % 关键词用“英文逗号”分割
  ckeywords={大语言模型推理, 错误识别, 强化学习},
  ekeywords={Large Language Model Reasoning, Error Identification, Reinforcement Learning},
}

\begin{cabstract}
近年来，大语言模型在问答系统、多模态交互、代码生成等诸多任务中取得了突破性进展，展现出较强的语言理解、知识整合与问题求解能力。
然而，大模型可能出现的错误制约其在实际场景中的应用。尤其是在需要多步推理的任务中，大语言模型不仅可能在最终答案上出现偏差，
还可能在推理链条的中间环节产生错误，推理中错误会逐步累积并放大，最终影响模型结果的可靠性与可解释性。
现有问答场景下的大语言模型错误检测与训练大多关注结果的正确性，而未对推理过程中可能产生的错误给予足够的重视，即使模型给出了看似正确的答案，
其推理过程仍可能存在问题，从而使模型产生可能的结果与过程不对齐，进而限制了大语言模型在高可靠性场景中的进一步应用。

因此，本文认为必须验证推理的严谨性。
具体而言，本文首先设计基于大语言模型的推理错误识别系统，对模型生成的中间推理步骤进行细粒度分析，
识别其中可能存在的事实性错误、逻辑性错误与推导一致性问题。在此基础上，进一步利用错误识别结果构建面向训练的增广数据，
使模型能够从典型错误样例及其修正过程中学习更加稳健的推理模式。最后，本文通过增广数据对模型进行优化训练，
以提升其在复杂问答任务中的推理准确性、过程一致性与结果可靠性。通过上述研究，
本文期望为大语言模型推理过程评估、错误诊断以及能力增强提供一种可行的方法路径。
使用本文所述方法构建的模型在多个基准测试上取得了比一般强化学习更好的表现。最后，本文基于错误识别
模型与回答模型的集成，构建了一个可迭代生成更可信回复的系统。

\end{cabstract}

\begin{eabstract}
In recent years, large language models have achieved breakthrough progress in numerous tasks such as question-answering systems, multimodal interaction, and code generation, demonstrating strong capabilities in language understanding, knowledge integration, and problem-solving.
However, the potential errors that large models may produce constrain their application in real-world scenarios. Particularly in tasks requiring multi-step reasoning, large language models may not only deviate in the final answer but also generate errors in the intermediate steps of the reasoning chain. These errors can accumulate and amplify gradually, ultimately affecting the reliability and interpretability of the model's results.
Current error detection and training of large language models in question-answering scenarios mainly focus on the correctness of the results, without paying sufficient attention to potential errors during the reasoning process. Even if the model provides seemingly correct answers, its reasoning process may still contain issues, leading to a misalignment between the model's output and the reasoning steps, which further limits the application of large language models in high-reliability scenarios.

Therefore, this paper argues that the rigor of reasoning must be verified. Specifically, this paper first designs a reasoning error identification system based on large language models, performing fine-grained analysis on the intermediate reasoning steps generated by the model to identify potential factual errors, logical errors, and inferential consistency issues. On this basis, the error identification results are further used to construct augmented data for training, enabling the model to learn more robust reasoning patterns from typical error examples and their correction processes. Finally, this paper optimizes the model through training on augmented data to enhance its reasoning accuracy, process consistency, and result reliability in complex question-answering tasks.
Through the above research, this paper aims to provide a feasible methodological path for evaluating the reasoning process of large language models, diagnosing errors, and enhancing their capabilities. Models constructed using the methods described in this paper have achieved better performance on multiple benchmark tests compared to general reinforcement learning. Finally, this paper builds an integrated system based on the error identification model and the answer model, creating a system capable of iteratively generating more credible responses.
\end{eabstract}
